\section*{Data and code availability}
\label{sec:availability}

No empirical data were collected, fitted or used in this work. The Lean
development, the finite models, the result snapshot and the scripts that check
them are in a public repository,
\url{https://github.com/ioannist/six-birds-transcendence}. The preprint's
archival DOI is \href{https://doi.org/\paperdoi}{doi:\paperdoi}.

The development builds with Lean 4.28.0 (\relLeanToolchain) using its core
library only. It has \relCstttModules{} modules and records
\relUniqueInventoryDeclarations{} distinct declarations in its inventory. It
proves \relProvedRows{} of the \relRegistryRows{} results it states and checks
the other \relCheckedCountermodels{} as finite countermodels, and
\relOpenRows{} are left open. The \relHeadlineDeclarations{} declarations we
cite depend only on the standard axioms \relHeadlineAxioms{}. Our
development imports no part of the vendored Six Birds library. That library
contains a non-standard axiom on which none of our declarations depends. The
finite models are specified in \relModelSpecs{} model files.

In the verification recorded on 2026-10-13 and 2026-10-14, a clean build ran \relLeanBuildJobs{} Lean
build jobs and succeeded, and the validator found no placeholder and no
untrusted axiom. All transcription checks between model files and Lean passed,
as did the validator's \relSelftestValidateLeanTestsPassed{} self-tests and
\relPaperTestsPassed{} paper tests. The repository also keeps the test suite of
an earlier Python scaffold that the Lean development superseded. There,
\relCstTtTestsPassed{} tests pass and \relCstTtTestsFailed{} fail, as they
already did when the Lean development was completed. Five of the failing
tests scan the superseded scaffold, and one tests that a historical file
builder is idempotent. No result here depends on the scaffold.

Every computed value in our figures and tables is read from the snapshot and
pinned to a Lean theorem that states it. The stage and callout numbers in the
schematic Figure~\ref{fig:three} are labels, not values. The kernel checks the Lean statements. That
those statements say what our prose says is not a kernel fact. We checked it in
a separate review, which set each formal statement in this paper beside the
printed Lean statement of every declaration it cites, together with that
declaration's axioms. The review used AI assistance, as disclosed below.

\paragraph{Disclosures.}
This preprint has not been peer reviewed. The repository's review records
document an internal review loop driven by the author. The Lean
formalization, the repository tooling, the preparation of the manuscript and
the statement-by-statement review used AI assistance. The author is responsible
for all claims.
